\chapter{高级排版技巧与模板进阶}
\label{chap:advanced}

本章将深入介绍 NUDT 论文模板的高级功能和使用技巧，包括复杂表格、交叉引用、符号列表、定制化设置等内容。

\section{复杂表格制作}
\label{sec:complex_tables}

\subsection{跨行跨列表格}
使用 \texttt{multirow} 和 \texttt{multicolumn} 创建复杂表格：

\begin{table}[htbp]
  \centering
  \caption{跨行跨列表格示例}
  \label{tab:complex_table}
  \begin{tabular}{|c|c|c|c|c|}
    \hline
    \multirow{2}{*}{\backslashbox{学期}{科目}} & \multicolumn{2}{c|}{上学期} & \multicolumn{2}{c|}{下学期} \\
    \cline{2-5}
    & 必修 & 选修 & 必修 & 选修 \\
    \hline
    \multirow{3}{*}{数学} & 高等数学 & 数学建模 & 概率统计 & 数值分析 \\
    & 线性代数 & 离散数学 & 复变函数 & 优化方法 \\
    & 数学分析 & 拓扑学 & 泛函分析 & 微分几何 \\
    \hline
    \multirow{2}{*}{物理} & 大学物理 & 近代物理 & 电动力学 & 量子力学 \\
    & 理论力学 & 固体物理 & 热力学 & 统计物理 \\
    \hline
  \end{tabular}
\end{table}

\subsection{长表格与续表}
使用 \texttt{longtable} 环境创建跨页表格：

\begin{longtable}{cccccc}
\caption{实验数据长表格示例}\label{tab:longtable_example}\\
\toprule[1.5pt]
{\hei 序号} & {\hei 日期} & {\hei 温度(℃)} & {\hei 湿度(\%)} & {\hei 压力(kPa)} & {\hei 备注} \\
\midrule[1pt]
\endfirsthead

\multicolumn{6}{c}{续表\thetable\hskip1em 实验数据长表格示例}\\

\toprule[1.5pt]
{\hei 序号} & {\hei 日期} & {\hei 温度(℃)} & {\hei 湿度(\%)} & {\hei 压力(kPa)} & {\hei 备注} \\
\midrule[1pt]
\endhead

\hline
\multicolumn{6}{r}{续下页}
\endfoot
\endlastfoot

1 & 2023-01-01 & 25.3 & 45.2 & 101.3 & 正常 \\
2 & 2023-01-02 & 24.8 & 46.1 & 101.2 & 正常 \\
3 & 2023-01-03 & 26.1 & 44.8 & 101.4 & 正常 \\
4 & 2023-01-04 & 25.7 & 45.5 & 101.3 & 正常 \\
5 & 2023-01-05 & 24.9 & 46.3 & 101.2 & 正常 \\
6 & 2023-01-06 & 25.2 & 45.8 & 101.3 & 正常 \\
7 & 2023-01-07 & 25.5 & 45.0 & 101.4 & 正常 \\
8 & 2023-01-08 & 24.7 & 46.5 & 101.2 & 正常 \\
9 & 2023-01-09 & 25.8 & 44.7 & 101.4 & 正常 \\
10 & 2023-01-10 & 25.1 & 45.9 & 101.3 & 正常 \\
11 & 2023-01-11 & 26.2 & 44.5 & 101.4 & 正常 \\
12 & 2023-01-12 & 25.6 & 45.3 & 101.3 & 正常 \\
13 & 2023-01-13 & 24.8 & 46.2 & 101.2 & 正常 \\
14 & 2023-01-14 & 25.9 & 44.6 & 101.4 & 正常 \\
15 & 2023-01-15 & 25.3 & 45.7 & 101.3 & 正常 \\
\bottomrule[1.5pt]
\end{longtable}

\subsection{可调宽表格}
自定义列类型实现自动换行和宽度调整：

\begin{table}[htbp]
  \centering
  \caption{可调宽表格示例}
  \label{tab:adjustable_table}
  \begin{tabularx}{0.9\linewidth}{|l|X|l|X|}
    \hline
    \textbf{编号} & \textbf{论文题目} & \textbf{作者} & \textbf{发表期刊/会议} \\
    \hline
    1 & 基于深度学习的图像识别算法研究 & 张三 & IEEE Transactions on Pattern Analysis and Machine Intelligence \\
    \hline
    2 & 量子计算在密码学中的应用 & 李四 & Nature Communications \\
    \hline
    3 & 高光谱遥感图像分类方法 & 王五 & Remote Sensing of Environment \\
    \hline
  \end{tabularx}
\end{table}

\section{交叉引用与超链接}
\label{sec:cross_ref}

\subsection{基本交叉引用}
模板自动处理各种交叉引用：

\begin{itemize}
  \item 章节引用：如第\ref{chap:intro}章介绍了基础用法
  \item 公式引用：公式\eqref{eq:newton}是牛顿第二定律
  \item 表格引用：表\ref{tab:basic_table}展示了基本表格样式
  \item 图片引用：图\ref{fig:single_image}是单张图片示例
  \item 定理引用：定理\ref{thm:pythagorean}是勾股定理
\end{itemize}

\subsection{超链接设置}
模板已配置超链接，在 PDF 中点击引用可跳转：

\begin{enumerate}
  \item 引用\cite{Black2001}可跳转到参考文献
  \item 点击"表\ref{tab:class_options}"可跳转到对应表格
  \item 点击"图\ref{fig:subfig_example}"可跳转到对应图片
\end{enumerate}

\section{符号列表与术语表}
\label{sec:symbols}

\subsection{使用 nomencl 宏包}
在导言区添加 \verb|\makenomenclature|，在正文中定义符号：

% 符号定义示例
\nomenclature[Aa]{$A$}{面积（Area）}{m$^2$}
\nomenclature[Av]{$v$}{速度（Velocity）}{m/s}
\nomenclature[At]{$t$}{时间（Time）}{s}
\nomenclature[Gr]{$\rho$}{密度（Density）}{kg/m$^3$}
\nomenclature[Ge]{$\epsilon$}{介电常数（Permittivity）}{F/m}
\nomenclature[Gm]{$\mu$}{磁导率（Permeability）}{H/m}

\subsection{使用 denotation 环境}
在 \texttt{data/denotation.tex} 中定义符号：

\begin{denotation}[3.5cm]
  \item[$c$] 光速 (Speed of light)，$3\times10^8$ m/s
  \item[$h$] 普朗克常数 (Planck constant)，$6.626\times10^{-34}$ J·s
  \item[$k_B$] 玻尔兹曼常数 (Boltzmann constant)，$1.381\times10^{-23}$ J/K
  \item[$e$] 元电荷 (Elementary charge)，$1.602\times10^{-19}$ C
  \item[$\pi$] 圆周率 (Pi)，3.1415926535
\end{denotation}

\section{定理类环境进阶}
\label{sec:theorems_adv}

\subsection{推论与命题}
\begin{corollary}[三角形内角和]
\label{cor:triangle_sum}
三角形内角和等于 $180^\circ$。
\end{corollary}

\begin{proposition}[连续函数性质]
\label{prop:continuous}
如果函数 $f(x)$ 在闭区间 $[a,b]$ 上连续，则：
\begin{enumerate}
  \item $f(x)$ 在 $[a,b]$ 上有界
  \item $f(x)$ 在 $[a,b]$ 上可取得最大值和最小值
  \item $f(x)$ 在 $[a,b]$ 上一致连续
\end{enumerate}
\end{proposition}

\subsection{例子与注释}
\begin{example}[导数计算]
\label{ex:derivative}
计算函数 $f(x) = x^2$ 在 $x=2$ 处的导数：
\[
f'(2) = \lim_{h \to 0} \frac{(2+h)^2 - 4}{h} = \lim_{h \to 0} (4 + h) = 4
\]
\end{example}

\begin{remark}
\label{rem:note}
注意：定理\ref{thm:pythagorean}仅适用于欧几里得空间中的直角三角形。在非欧几何中，该定理不成立。
\end{remark}

\section{算法环境进阶}
\label{sec:algorithms_adv}

\subsection{复杂算法示例}
\begin{algorithm}[htbp]
  \caption{支持向量机训练算法}
  \label{alg:svm}
  \begin{algorithmic}[1]
    \REQUIRE 训练集 $D = \{(\mathbf{x}_i, y_i)\}_{i=1}^n$，正则化参数 $C$，核函数 $K$
    \ENSURE 支持向量 $\mathbf{SV}$，拉格朗日乘子 $\boldsymbol{\alpha}$
    
    \STATE 初始化 $\boldsymbol{\alpha} \leftarrow \mathbf{0}$，$b \leftarrow 0$
    \STATE 计算核矩阵 $K_{ij} = K(\mathbf{x}_i, \mathbf{x}_j)$
    
    \WHILE{未收敛}
      \FOR{$i = 1$ to $n$}
        \STATE 计算预测值：$f_i = \sum_{j=1}^n \alpha_j y_j K_{ij} + b$
        \STATE 计算误差：$E_i = f_i - y_i$
        
        \IF{$(\alpha_i < C$ and $y_i E_i < -\epsilon)$ or $(\alpha_i > 0$ and $y_i E_i > \epsilon)$}
          \STATE 随机选择 $j \neq i$
          \STATE 计算边界：$L = \max(0, \alpha_j - \alpha_i)$，$H = \min(C, C + \alpha_j - \alpha_i)$
          
          \IF{$|L - H| < \epsilon$}
            \STATE \textbf{continue}
          \ENDIF
          
          \STATE 更新 $\alpha_j$：$\alpha_j^{\text{new}} = \alpha_j - \frac{y_j(E_i - E_j)}{\eta}$
          \STATE 裁剪：$\alpha_j^{\text{new}} = \min(\max(\alpha_j^{\text{new}}, L), H)$
          \STATE 更新 $\alpha_i$：$\alpha_i^{\text{new}} = \alpha_i + y_i y_j (\alpha_j - \alpha_j^{\text{new}})$
          
          \STATE 更新阈值 $b$
        \ENDIF
      \ENDFOR
    \ENDWHILE
    
    \STATE 提取支持向量：$\mathbf{SV} = \{\mathbf{x}_i | \alpha_i > 0\}$
    \RETURN $\boldsymbol{\alpha}, b, \mathbf{SV}$
  \end{algorithmic}
\end{algorithm}

\subsection{跨页算法}
使用自定义的 \texttt{breakablealgorithm} 环境：

\begin{breakablealgorithm}
  \caption{跨页算法示例}
  \begin{algorithmic}[1]
    \REQUIRE 数据集 $X \in \mathbb{R}^{n \times d}$，聚类数 $k$
    \ENSURE 聚类中心 $C$，标签 $L$
    
    \STATE 随机初始化聚类中心 $C = \{\mathbf{c}_1, \ldots, \mathbf{c}_k\}$
    \REPEAT
      \STATE 分配步骤：为每个样本分配最近聚类中心
      \FOR{$i = 1$ to $n$}
        \STATE $l_i = \arg\min_{j} \|\mathbf{x}_i - \mathbf{c}_j\|^2$
      \ENDFOR
      
      \STATE 更新步骤：重新计算聚类中心
      \FOR{$j = 1$ to $k$}
        \STATE $\mathbf{c}_j = \frac{1}{|S_j|} \sum_{\mathbf{x} \in S_j} \mathbf{x}$
        \COMMENT{$S_j = \{\mathbf{x}_i | l_i = j\}$}
      \ENDFOR
      
      \STATE 计算目标函数：$J = \sum_{i=1}^n \|\mathbf{x}_i - \mathbf{c}_{l_i}\|^2$
    \UNTIL{$J$ 收敛或达到最大迭代次数}
    
    \STATE 这是算法的第一部分，如果内容足够多，会自动跨页显示
    \STATE 继续添加一些步骤...
    \FOR{$i = 1$ to $10$}
      \STATE 步骤 $i$：处理数据
      \IF{$i$ 是偶数}
        \STATE 偶数步骤特殊处理
      \ELSE
        \STATE 奇数步骤正常处理
      \ENDIF
    \ENDFOR
    
    \STATE 算法结束，返回结果
    \RETURN $C, L$
  \end{algorithmic}
\end{breakablealgorithm}

\section{代码高亮进阶}
\label{sec:code_adv}

\subsection{C++ 代码示例}
\begin{lstlisting}[language=C++, caption=C++矩阵运算示例]
#include <iostream>
#include <vector>
#include <cmath>

class Matrix {
private:
    std::vector<std::vector<double>> data;
    int rows, cols;
    
public:
    Matrix(int r, int c) : rows(r), cols(c) {
        data.resize(r, std::vector<double>(c, 0.0));
    }
    
    // 矩阵乘法
    Matrix multiply(const Matrix& other) const {
        if (cols != other.rows) {
            throw std::invalid_argument("矩阵维度不匹配");
        }
        
        Matrix result(rows, other.cols);
        for (int i = 0; i < rows; ++i) {
            for (int j = 0; j < other.cols; ++j) {
                double sum = 0.0;
                for (int k = 0; k < cols; ++k) {
                    sum += data[i][k] * other.data[k][j];
                }
                result.data[i][j] = sum;
            }
        }
        return result;
    }
    
    // 矩阵转置
    Matrix transpose() const {
        Matrix result(cols, rows);
        for (int i = 0; i < rows; ++i) {
            for (int j = 0; j < cols; ++j) {
                result.data[j][i] = data[i][j];
            }
        }
        return result;
    }
    
    // 打印矩阵
    void print() const {
        for (int i = 0; i < rows; ++i) {
            for (int j = 0; j < cols; ++j) {
                std::cout << data[i][j] << " ";
            }
            std::cout << std::endl;
        }
    }
};

int main() {
    Matrix A(2, 3);
    Matrix B(3, 2);
    // 初始化矩阵...
    Matrix C = A.multiply(B);
    C.print();
    return 0;
}
\end{lstlisting}

\subsection{LaTeX 代码示例}
\begin{lstlisting}[language=TeX, caption=LaTeX自定义命令示例]
% 在 mynudt.sty 中添加自定义命令
% 数学符号快捷命令
\newcommand{\R}{\mathbb{R}} % 实数集
\newcommand{\C}{\mathbb{C}} % 复数集
\newcommand{\Z}{\mathbb{Z}} % 整数集
\newcommand{\N}{\mathbb{N}} % 自然数集

% 微分算子
\newcommand{\diff}{\mathrm{d}}
\newcommand{\pdiff}[2]{\frac{\partial #1}{\partial #2}}

% 期望和方差
\newcommand{\E}{\mathbb{E}}
\newcommand{\Var}{\mathrm{Var}}

% 自定义定理环境
\newtheorem{assumption}{假设}[chapter]
\newtheorem{property}{性质}[section]

% 算法中的特殊命令
\algnewcommand{\algorithmicforeach}{\textbf{for each}}
\algnewcommand{\ForEach}[1]{\algorithmicforeach\ #1\ \algorithmicdo}

% 自定义列表环境
\newenvironment{mylist}{
  \begin{enumerate}[label=(\arabic*), leftmargin=*]
}{
  \end{enumerate}
}
\end{lstlisting}

\section{参考文献管理进阶}
\label{sec:bib_adv}

\subsection{BibLaTeX 高级功能}
使用 BibLaTeX 的强大功能：

\begin{itemize}
  \item 按章节拆分参考文献
  \item 多文献库合并
  \item 分类显示文献（期刊、会议、专著等）
  \item 自定义引用样式
\end{itemize}

\subsection{参考文献分类显示}
\begin{lstlisting}[language=TeX, caption=分类显示参考文献]
% 在导言区添加
\DeclareBibliographyCategory{journal}
\DeclareBibliographyCategory{conference}
\DeclareBibliographyCategory{book}

% 在正文中添加文献到分类
\addtocategory{journal}{Black2001,Bauters2011}
\addtocategory{conference}{Akahane2003}
\addtocategory{book}{Coldren2012,Agrawal2000}

% 分类打印参考文献
\printbibliography[category=journal, title=期刊论文]
\printbibliography[category=conference, title=会议论文]
\printbibliography[category=book, title=专著]
\end{lstlisting}

\section{模板定制与扩展}
\label{sec:customization}

\subsection{修改页面布局}
在导言区调整页面边距：

\begin{lstlisting}[language=TeX]
% 调整页边距
\geometry{
  top=25mm,      % 上边距
  bottom=30mm,   % 下边距
  left=35mm,     % 左边距
  right=25mm,    % 右边距
  headheight=10mm,
  headsep=2mm,
  footskip=10mm
}
\end{lstlisting}

\subsection{自定义章节格式}
修改章节标题格式：

\begin{lstlisting}[language=TeX]
% 修改章节格式
\titleformat{\chapter}
  {\filcenter\heiti\sanhao[1.5]}  % 格式
  {第\zhnumber{\thechapter}章}    % 标签
  {1em}                          % 间距
  {}                             % 标题前代码
  [\titlerule]                   % 标题后代码
  
% 修改小节格式
\titleformat{\section}
  {\filcenter\bfseries\heiti\sihao[1.3]}
  {\thesection}
  {1em}
  {\underline}
\end{lstlisting}

\subsection{添加水印}
添加论文水印：

\begin{lstlisting}[language=TeX]
% 需要添加包
\usepackage{draftwatermark}
\SetWatermarkText{国防科技大学}
\SetWatermarkScale{0.8}
\SetWatermarkColor[gray]{0.9}
\SetWatermarkAngle{45}
\end{lstlisting}

\section{编译与调试}
\label{sec:compilation}

\subsection{编译流程}
推荐编译流程：

\begin{enumerate}
  \item \verb|xelatex mainpaper| - 第一次编译，生成辅助文件
  \item \verb|biber mainpaper|   - 处理参考文献
  \item \verb|xelatex mainpaper| - 第二次编译，更新引用
  \item \verb|xelatex mainpaper| - 第三次编译，确保所有引用正确
\end{enumerate}

\subsection{常见问题解决}
\begin{table}[htbp]
  \centering
  \caption{常见编译问题及解决方案}
  \label{tab:compile_issues}
  \begin{tabularx}{\linewidth}{lX}
    \toprule[1.5pt]
    {\hei 问题} & {\hei 解决方案} \\
    \midrule[1pt]
    参考文献引用显示为 [?] & 运行 \texttt{biber} 命令处理参考文献 \\
    图片引用错误 & 检查标签是否重复，运行多次编译 \\
    表格/图片位置不合适 & 调整 \texttt{[htbp]} 参数，使用 \texttt{H} 强制位置 \\
    数学公式编号错误 & 检查 \texttt{equation} 环境是否嵌套 \\
    字体缺失 & 检查字体选项设置，确保系统安装相应字体 \\
    \bottomrule[1.5pt]
  \end{tabularx}
\end{table}

\section{本章小结}
本章详细介绍了 NUDT 论文模板的高级功能：
\begin{itemize}
  \item 复杂表格的制作与排版技巧
  \item 交叉引用与超链接的使用
  \item 符号列表和术语表的创建
  \item 定理、算法环境的进阶应用
  \item 代码高亮和参考文献管理
  \item 模板定制与扩展方法
  \item 编译流程和问题调试
\end{itemize}

通过掌握这些高级技巧，您可以创建出专业、美观的学术论文，满足各种排版需求。